Este trabalho delimitou o problema de prompt injection indireta em um agente que processa documentos de Contas a Pagar e atua sobre Odoo. A arquitetura proposta separa interpretação, autorização e execução: o LLM produz uma proposta tipada; regras determinísticas confrontam a proposta com pedido, fornecedor, beneficiário e estado confiáveis; ações críticas exigem aprovação vinculada; e o executor opera com menor privilégio, idempotência e auditoria. C0 e C1 serão comparados sob o mesmo modelo, documento, prompt-base, schema e fixture.

Na etapa registrada por este artigo, foi demonstrada a viabilidade técnica de partes do aparato experimental. Odoo 19.0 Community foi fixado em uma revisão reproduzível; M1/M2 comprovaram o fluxo PDF--proposta--rascunho e o reset de uma fixture sintética. Incrementos posteriores acrescentaram contrato comum v2, caminhos privados C0/C1, captura independente, restauração de banco e filestore e registro candidato que preserva falhas. Essas provas reduzem riscos de instrumentação, mas não equivalem a um runner científico integrado e revisado.

O piloto de propostas, executado com o modelo real, trouxe o primeiro resultado empírico. Em duas rodadas, a segunda com estratégias de ataque revisadas, nenhuma das 72 propostas adversariais de C0 incorporou o objetivo do atacante, e as propostas benignas foram exatas em 36 de 36 execuções em C0 e em 35 de 36 em C1. Com esse modelo, esse prompt-base e ataques estáticos, a redução do sucesso dos ataques por C1 não pode ser demonstrada, porque a baseline já não cede. Isso não prova robustez do modelo nem dispensa a fronteira determinística, cuja função é independer do que o modelo proponha.

Esses resultados não respondem ainda se C1 reduz ASR-state ou UACR, se preserva BTSR no estado do ERP ou qual sobrecarga integral introduz. A coleta confirmatória não ocorreu, e os testes privados não abrem o método público de execução. Assim, a conclusão permitida é restrita: existem componentes observáveis e restauráveis, e o modelo resistiu aos ataques do piloto, mas a eficácia e o custo da proteção permanecem questões abertas.

As próximas etapas são decidir o enquadramento estimativo proposto, ligar o ciclo integrado aos pares do piloto e medir o estado, concluir a revisão, congelar o protocolo e executar os pares randomizados. O núcleo preserva 144 execuções piloto e 1.080 principais, com três repetições, totalizando 1.224 sem retries. O artigo final deverá responder \hyperlink{rq1}{RQ1}--\hyperlink{rq3}{RQ3} com denominadores, efeitos, intervalos e descrição de falhas, sem hipótese de aceitabilidade de custo. Segundo modelo, ablação e ataques adaptados à defesa são trabalhos futuros. Sua ausência limita a generalização entre modelos, a atribuição causal por camada e a resistência adaptativa. Pagamentos, documentos reais, OCR e comportamento de revisores humanos continuarão fora do escopo.

A contribuição científica dependerá da combinação entre arquitetura executável e avaliação honesta. Se C1 reduzir violações sem ultrapassar a margem de utilidade, o estudo fornecerá evidência situada para integrações semelhantes. Se não reduzir, ou se o custo for excessivo, o resultado continuará relevante ao mostrar os limites práticos das camadas escolhidas. Em ambos os casos, a conclusão final será determinada pelos dados congelados, não pela expectativa de que a arquitetura funcione.
